\documentclass[a4paper,11pt]{article}
\usepackage[utf8]{inputenc}
\usepackage[T1]{fontenc}
\usepackage[italian]{babel}
\usepackage[margin=2.5cm]{geometry}
\usepackage{titlesec}
\usepackage{parskip}
\usepackage{xcolor}
\usepackage{amssymb}
\usepackage{amsmath}
\usepackage{algorithm}
\usepackage{algpseudocode}

\definecolor{grigio}{RGB}{100,100,100}

\titleformat{\section}{\large\bfseries}{}{0em}{}
\titlespacing{\section}{0pt}{1.8em}{0.4em}

\newcommand{\domanda}[2]{%
  \vspace{0.6em}
  \noindent\textbf{#1}\\[0.2em]
  \noindent#2
  \vspace{1em}
  \par
}

\begin{document}

\begin{center}
  {\LARGE\bfseries Ottimizzazione Combinatoria}\\[0.5em]
  {\large Domande di teoria 2025--26}
\end{center}

\vspace{1.5em}

\section{1 \quad Ottimizzazione lineare}

\domanda{1.1 \quad Descrivere i quattro casi di soluzione di un problema di ottimizzazione.}{
\begin{enumerate}
    \item Soluzione ottima unica: unico punto che massimizza la funzione obbiettivo rispetto ai vincoli
    \item Soluzione ottima infinite: infiniti punti che massimizzano la funzione obbiettivo rispetto ai vincoli
    \item Problema illimitato: la soluzione tende all'infinito
    \item Problema inamissibile: i vincoli non permettono nessun valore 
\end{enumerate}
}

\domanda{1.2 \quad Definire problema di ottimizzazione, di esistenza e di certificato.}{
Un problema di ottimizzazione e' definito da:
\begin{itemize}
    \item $S$ come insieme soluzioni ammissibili
    \item $f \, : \, S \to \mathbb{R}$ come funzione obbiettivo
    \item $\max$ (o $\min$)
\end{itemize}
t.c. $f(x^*) = \max f$\\
\vspace{1px}\\
Un problema di esistenza: \textit{Esiste almeno una soluzione valida?}
\begin{align*}
    S \neq \emptyset
\end{align*}
\vspace{1px}\\
Un problema di certificato indica la dimostrazione del risultato che abbiamo ottenuto:
\begin{itemize}
    \item Se esiste una soluzione controlliamo in punto $x \in S$ se rispetta i vincoli
    \item Se non esiste una soluzione allora dobbiamo trovare dei vincoli che si contraddicono!
\end{itemize}
}

\domanda{1.3 \quad Definire un algoritmo di tipo primale e duale e fornirne un esempio.}{
\begin{itemize}
    \item Primale: parte da una soluzione ammissibile per il primale e la ottimizza (mantenendo l'ammissibile)
    \item Duale: parte della soluzione ottimale per il duale (non ammissibile per il primale) e raggiunge l'ammissibilita' (mantenendo l'ottimo)
\end{itemize}
Un esempio e' Simplex e Simplex duale.
}

\domanda{1.4 \quad Definire un rilassamento di problema di ottimizzazione e illustrarne un esempio.}{
Dato un problema di ottimizzazione $P=(s, f \max)$, un problema di rilassamento e' un problema di ottimizzazione $P'=(S', f', \max)$ t.c.
\begin{enumerate}
    \item $S \subseteq S'$
    \item $f'(x) \geq f(x)$ 
\end{enumerate}
$\Rightarrow$ l'ottimo di $P$ e' minore uguale al ottimo di $P'$. abbiamo un upper-bound su $P$.\\
\vspace{1px}\\
Un esempio e'
\begin{align*}
    \max \quad &x \\
    \text{s.t. } A&x\leq b \quad x \in \mathbb{N}\cup\{0\}
\end{align*}
Il rilassamento $P'$ e':
\begin{align*}
    \max \quad &x \\
    \text{s.t. } A &x\leq b \quad x \geq 0
\end{align*}
Quindi $x$ puo' essere anche $1.5, 3.2, \dots$, ovvero numeri non interi.
}

\domanda{1.5 \quad Discutere ed illustrare con un esempio la modellazione con vincoli lineari di relazioni logiche (almeno un caso).}{
Per rappresentare le relazioni logiche in un problema di ottimizzazione utilizziamo una formulazione specifica. \\
\vspace{1px}\\
Tipo per descrivere il problema di massimo sulla base $0 \leq x \leq 10$ con i vincoli $x \leq 3 \lor x \geq 7$:\\
\texttt{Variabili}
\begin{align*}
    &y = \begin{cases}
        0 \text{ se } x \leq 3\\
        1 \text{ se } x \geq 7
    \end{cases}\\
    &M=10
\end{align*}
\texttt{Funzione obbiettivo}
\begin{align*}
    \max \quad &x \\
    \text{s.t } \quad &x \geq 0\\
    &x \leq 10
\end{align*}
\texttt{Vincoli}
\begin{align*}
    x \leq 3 + My \\
    x \geq 7 - M(1-y)
\end{align*}
\vspace{1px}\\
Per \texttt{and} basta scrivere i due vincoli, mentre per \texttt{implicazione} su $A \to B$ dobbiamo fare
\begin{align*}
    y_A \leq y_B
\end{align*}
ovvero se $A$ e' vero allora $B$ lo e', mentre se $A$ non e' vero allora $B$ puo' essere qualsiasi cosa.
}

\domanda{1.6 \quad Illustrare la modellazione con vincoli lineari di funzioni lineari a tratti e convesse.}{
Dato una funzione del tipo
\begin{align*}
    f(x) = \max(a_1x+b_1,\ a_2x+b_2,\ \dots,\ a_kx+b_k)
\end{align*}
Abbiamo una funzione convessa composta da rette.\\
Se volessimo trovare il minimo di essa?
\begin{align*}
    &\min \quad f(x)\\
    =&\min \quad \max(a_1x+b_1,\ a_2x+b_2,\ \dots,\ a_kx+b_k)
\end{align*}
Il problema e' che $\max$ non e' lineare! Introduciamo quindi una variabile ausiliaria $z$:
\begin{align*}
    \min \quad &z\\
    \text{s.t. }\quad &z \geq a_1x+b_1\\
    &z \geq a_2x+b_2\\
    &\vdots\\
    &z \geq a_kx+b_k
\end{align*}
Questo porta precisamente che $z = \max(a_1x+b_1,\ a_2x+b_2,\ \dots,\ a_kx+b_k)$.
}

\domanda{1.7 \quad Definire una soluzione base ed una soluzione base ammissibile di un problema PL.}{
Dato un problema di ottimizzazione in forma standard
\begin{align*}
    \min \quad &x\\
    \text{s.t.} \quad &Ax = b, \quad x \geq 0
\end{align*}
scegliamo $m$ colonne linearmente indipendenti di $A$ formando la \textbf{base} $A_B$, con $A_N$ le colonne fuori base.\\
Ponendo $x_N = 0$ otteniamo la \textbf{soluzione di base}:
\begin{align*}
    \hat{x}_B = A^{-1}_B b
\end{align*}
Se vale anche $\hat{x}_B \geq 0$, allora $\hat{x}$ è una \textbf{soluzione di base ammissibile} (SBA).
}

\domanda{1.8 \quad Enunciare il teorema di Motzkin (decomposizione) per un poliedro.}{
Ogni poliedro $P=Ax\leq b$ puo' essere riscritto come:
\begin{align*}
    P=Q+C
\end{align*}
\begin{itemize}
    \item $Q$ guscio convesso di vertici $v_i$
    \item $C$ cono di direzioni $d_j$
\end{itemize}
Ogni punto puo' essere scritto come
\begin{align*}
    x = \sum_i \lambda_i v_i + \sum_j \mu_j d_j
\end{align*}
}

\domanda{1.9 \quad Enunciare e dimostrare la proprietà fondamentale della PL.}{
\texttt{Enunciato}: se un problema di PL ha soluzione ottima allora il suo poliedro ha soluzione ottima in uno dei suoi vertici!\\
\vspace{1px}\\
\texttt{Dimostrazione}: Dal teorema di Motzkin sappiamo che abbiamo vertici e direzioni
\begin{enumerate}
    \item Essendo $x^*$ massimo non abbiamo direzioni, ovvero il cono non contribuisce
    \item \begin{align*}
        z^* = x^* = &\sum_i \lambda_i v_i + \sum_j \mu_j d_j \leq \sum_i \lambda_i v^*\\
        &\sum_i \lambda_i v_i \leq \sum_i \lambda_i v^*=v^*\\
        \Rightarrow & z^* \leq v^*
    \end{align*}
    ma dato che $z^*$ e' soluzione ottima allora
    \begin{align*}
        z^* = v^*
    \end{align*}
    Quindi $v_i$ e' ottimo!
\end{enumerate}
}

\domanda{1.10 \quad Definire il problema primale e duale e fornire un esempio.}{
Dato un problema primale
\begin{align*}
    \min \quad &c^Tx \\
    \text{s.t. }\quad &Ax \geq b \quad x\geq 0
\end{align*}
Possiamo definire il duale come
\begin{align*}
    \max \quad &b^T y\\
    \text{s.t.}\quad &A^Ty \leq c \quad y \geq 0
\end{align*}
\vspace{1px}\\
Un esempio:
\begin{align*}
    \min \quad &3x_1 + 2x_2 \\
    \text{s.t.}\quad &x_1+x_2 \geq 4\\
    &2x_1+x_2 \geq 6\\
    &x_1,x_2 \geq 0
\end{align*}
Il duale e':
\begin{align*}
    \max \quad & 4y_1 + 6y_2\\
    &y_1 + 2y_2 \leq 3\\
    &y_1 + y_2 \leq 2
\end{align*}
L'ottimo coindice, il vantaggio e' che il duale potrebbe essere piu' facile da calcolare rispetto il primale.
}

\domanda{1.11 \quad Discutere l'interpretazione delle variabili duali.}{
Una variabile $i$ duale indica il valore marginale del vincolo
\begin{align*}
    y_i = \frac{\partial z^*}{\partial b_i}
\end{align*}
\begin{itemize}
    \item $y_i=0$ avere piu' risorse $i$ non migliora verso ottimo
    \item $y_i >0$ avere di piu' ci porta a ottimo, con un incremento di $y_i$
\end{itemize}
}

\domanda{1.12 \quad Enunciare e dimostrare il teorema debole della dualità.}{
\texttt{Enunciato:} Per ogni soluzione ammissibile $x$ del primale, per ogni soluzione ammisibile $y$ del duale:
\begin{align*}
    b^Ty \leq c^T x
\end{align*}
\vspace{1px}\\
\texttt{Dimostrazione}:
\begin{align*}
    1. \quad &Ax \geq b\\
    &y^TAx \geq y^T b \\
    &y^TAx \geq b^T y\\
    2. \quad &A^T y \leq c\\
    &x^T A^T y \leq x^T c\\
    &x^T A^T y \leq c^T x\\
\end{align*}
Sappiamo che $y^TAx = (y^TAx)^T = x^T A^T y$
\begin{align*}
    \Rightarrow &y^TAx \leq c^T x\\
\end{align*}
Allora
\begin{align*}
    \begin{cases}
    &y^TAx \geq b^Ty \\
    &y^TAx \leq c^T x
    \end{cases}\\
    \Rightarrow b^Ty \leq c^Tx
\end{align*}
}

\domanda{1.13 \quad Definire una direzione ammissibile e di crescita.}{
Data una \texttt{soluzioni ammissibile} $x \in P$, definiamo direzione ammissibile
\begin{align*}
    x + \epsilon d \in P \quad \exists \epsilon > 0
\end{align*}
Un vettore $d$ lo definiamo \texttt{direzione di crescita} se 
\begin{align*}
    c^T d > 0
\end{align*}
Cioè muovendosi nella direzione $d$ il valore della funzione obiettivo aumenta.
}

\domanda{1.14 \quad Enunciare la condizione di ottimalità per la PL.}{
Dato un problema di PL
\begin{align*}
    \min \quad c^T &x \\
    \text{s.t. } \quad A&x=b \quad x \geq 0
\end{align*}
\texttt{Enunciato:} Sia $x^*$ soluzione ammissibile alla base $B$, allora:
\begin{align*}
    \hat{c}_N^T=c_N^T- \hat{y}^T A_N \geq 0 
\end{align*}
}

\section{2 \quad Ottimizzazione su rete}

\domanda{2.1 \quad Definire il problema del flusso massimo su una rete ed il relativo modello matematico.}{
Una rete di flusso e' una quadrupla $G=(V, A, s, t)$ 
\begin{itemize}
    \item $V$ insieme dei nodi
    \item $S$ insieme degli archi orientati con capacita' $c_{ij}$
    \item $s$ nodo sorgente
    \item $t$ nodo pozzo
\end{itemize}
Il modello matematico si formula come 
\begin{align*}
    \max \quad & v = \sum_{j \,:\, (s,j) \in A} x_{sj} - \sum_{j \,:\, (j,s) \in A} x_{js} \\[6pt]
    & \sum_{j \,:\, (j,i) \in A} x_{ji} - \sum_{j \,:\, (i,j) \in A} x_{ij} = 0
      \qquad \forall\, i \in V \setminus \{s, t\} \\[4pt]
    & 0 \leq x_{ij} \leq c_{ij} \qquad \forall\, (i,j) \in A
\end{align*}
Con $x_{ij}$ flusso su arco $(i,j) \in A$.
}

\domanda{2.2 \quad Definire il problema del flusso di costo minimo su una rete ed il relativo modello matematico.}{
La rete e' una quintupla $G=(V, A, b_i, c_{ij}, j_{ij})$
\begin{itemize}
    \item $V$ insieme dei nodi
    \item $A$ insieme degli archi
    \item $b_i$ bilancio del nodo $i$
    \item $c_{ij}$ capacita' dell'arco $(i,j)$
    \item $k_{ij}$ costo unitario del flusso
\end{itemize}
Il modello matematico lo modelliamo come
\begin{align*}
    &\min \quad \sum_{(i,j) \in A} k_{ij}x_{ij}
\end{align*}
soggetto a
\begin{align*}
    & \sum_{j \,:\, (i,j) \in A} x_{ij} - \sum_{j \,:\, (j,i) \in A} x_{ji} = b_i \qquad \forall i \in V\\
    &0 \leq x_{ij} \leq c_{ij} \qquad \forall (i,j) \in A
\end{align*}
}

\domanda{2.3 \quad Data una rete, definire un taglio s-t ed illustrarne le proprietà.}{
Data una rete $G$, definiamo taglio $s-t$ una tupla ($S, T$) t.c.
\begin{itemize}
    \item $s \in S$
    \item $t \in T$
    \item $T = V \backslash S$
\end{itemize}
La capacita' di taglio e' definita 
\begin{align*}
    C(S, T) = \sum_{(i,j)\in S} c_{ij}
\end{align*}
Le proprieta' sono
\begin{itemize}
    \item Il flusso netto attraverso ogni taglio è uguale al valore del flusso
    \item Ogni taglio fornisce un limite superiore al valore del flusso
    \item Teorema Max-Flow Min-Cut 
\end{itemize}
}

\domanda{2.4 \quad Data una rete ed un flusso, definire il grafo residuo corrispondente.}{
\textit{Da svolgere nel esercizio dato dal prof.}
}

\domanda{2.5 \quad Definire un cammino aumentante per il flusso massimo sul grafo residuo.}{
Un cammino aumentante e' un cammino valido da $s$ a $t$ nel grafo residuo.
}

\domanda{2.6 \quad Discutere la correttezza dell'algoritmo di Ford-Fulkerson (ev.\ solo enunciati o specifiche fasi).}{
\begin{algorithm}
\caption{Ford-Fulkerson}
\begin{algorithmic}[H]
\State $x_{ij} = 0$ \textbf{per ogni} $(i,j) \in A$
\While{esiste un cammino aumentante $P$ in $G_x$}
    \State $\delta = \min \quad \text{residuo}_{ij}$
    \State aggiorna $x$ lungo $P$ con bottleneck $\delta$
\EndWhile\\
\Return $x$
\end{algorithmic}
\end{algorithm}\\
Correttezza:
\begin{enumerate}
    \item Il flusso resta ammissibile
    \item Quando si ferma, il flusso è massimo
    \item Termina sempre
\end{enumerate}
}

\domanda{2.7 \quad Enunciare e dimostrare il teorema max-flow-min-cut per il flusso massimo.}{
\texttt{Enunciato}: Per qualsiasi rete di flusso le tre condizioni sono equivalenti:
\begin{enumerate}
    \item $x$ è flusso massimo
    \item Non esiste nessun cammino aumentante in $G_x$
    \item Esiste un taglio $(S, \bar{S})$ t.c. $v(x) = C(S, \bar{S})$
\end{enumerate}

\texttt{Dimostrazione} ($1 \Rightarrow 2 \Rightarrow 3 \Rightarrow 1$):

\textbf{($1 \Rightarrow 2$)} Se esistesse un cammino aumentante in $G_x$,
si potrebbe aumentare $v(x)$, contraddicendo la massimalità. $\square$

\textbf{($2 \Rightarrow 3$)} Sia $S = \{\text{nodi raggiungibili da } s
\text{ in } G_x\}$. Poiché $t \notin S$, la coppia $(S, \bar{S})$ è un
taglio s-t. Per costruzione, ogni arco diretto del taglio è saturo
($x_{ij} = c_{ij}$) e ogni arco inverso è vuoto ($x_{ij} = 0$).
Dalla Proprietà 1 sui tagli:
\[
    v(x) = \sum_{A^+(S)} x_{ij} - \sum_{A^-(S)} x_{ij}
         = \sum_{A^+(S)} c_{ij} = C(S, \bar{S}) \qquad \square
\]

\textbf{($3 \Rightarrow 1$)} Dalla Proprietà 2 sappiamo che per qualsiasi
flusso $x'$ e qualsiasi taglio: $v(x') \leq C(S, \bar{S}) = v(x)$.
Quindi $x$ è massimo. $\square$

\textbf{Corollario}: combinando $1 \Rightarrow 2 \Rightarrow 3$ si ottiene
che quando Ford-Fulkerson termina, il taglio $(S, \bar{S})$ costruito è
minimo e il flusso corrente è massimo, quindi:
\[
    \max v(x) = \min C(S, \bar{S})
\]
}

\domanda{2.8 \quad Definire uno pseudoflusso e discutere le condizioni della sua ammissibilità.}{
Data una rete $G$ definiamo \texttt{pseudoflusso} una funzione che soddifsfi solo 
\begin{align*}
    0 \leq x_{ij} \leq c_{ij} \qquad \forall (i,j) \in A
\end{align*}
(La conservazione del flusso non è richiesta).\\
\vspace{1px}\\
Uno pseudoflusso e' flusso ammissibile se e solo se tutti i nodi sono sbilanciati
\begin{align*}
    e_i(x)=0 \qquad \forall i \in V
\end{align*}
}

\domanda{2.9 \quad Definire un cammino aumentante per il flusso a costo minimo sul grafo residuo.}{
Un cammino aumentante per il flusso a costo minimo e' un camminimo di costo minimo nel grafo residuo $G_x$ da un nodo di surplus e un nodo di decifit.
}

\domanda{2.10 \quad Definire uno pseudoflusso minimale.}{
Uno pseudoflusso e' minimale se nel grafo residuo $G_x$ non esistono cicli di costo negativo.
}

\domanda{2.11 \quad Discutere l'effetto di un cammino aumentante a costo minimo su uno pseudoflusso minimale.}{
\texttt{Teorema}: Se $x$ e' minimale e $P$ e' un camminimo di costo minimo in $G_x$ allora lo pseudoflusso $x'$ ottenuto aumentando lungo $P$ e' ancora minimale.
}

\domanda{2.12 \quad Discutere come ottenere un flusso ammissibile, se esiste, per un problema di flusso a costo minimo.}{
Si procede in due fasi: si costruisce un pseudoflusso minimale iniziale, poi lo si rende ammissibile tramite augmentation successive, controllando se la feasibility è raggiungibile. Cosi' otteniamo un flusso ammissibile e ottimo.
}

\end{document}